% Official Solution (Problem Sheet 11, Exercise 1)
\begin{exercisesolution}[Solution to \cref{exc:extending_linear_map_by_zero}]
\label{sol:extending_linear_map_by_zero}
The extended map $T \colon V \to W$ is \textbf{not} linear in general.

\textbf{Proof:}
Because $S$ is non-trivial, there exists a vector $u \in U$ such that $S(u) \ne 0_W$. Since $U \subsetneq V$ is a proper subspace, we can choose a vector $v_0 \in V \setminus U$.

Consider the sum $u + v_0 \in V$. We claim that $u + v_0 \notin U$: if $u + v_0 \in U$, then $v_0 = (u + v_0) - u \in U$ because $U$ is a subspace, contradicting the choice of $v_0 \in V \setminus U$.

Now evaluate $T$ on $u$, $v_0$, and $u + v_0$:
\begin{itemize}
    \item $u \in U \implies T(u) = S(u) \ne 0_W$;
    \item $v_0 \in V \setminus U \implies T(v_0) = 0_W$;
    \item $u + v_0 \in V \setminus U \implies T(u + v_0) = 0_W$.
\end{itemize}
Consequently,
\[
T(u + v_0) = 0_W \ne S(u) = T(u) + T(v_0).
\]
Hence $T$ fails additivity and is therefore not a linear map.
\exinfo{Official Solution (Problem Sheet 11, Exercise 1). Proves non-linearity of extending a non-trivial map by zero outside a proper subspace.}
\end{exercisesolution}

% Official Solution (Problem Sheet 11, Exercise 2)
\begin{exercisesolution}[Solution to \cref{exc:additive_maps_over_rationals}]
\label{sol:additive_maps_over_rationals}
Let $f \colon V \to W$ be an additive map, so $f(x + y) = f(x) + f(y)$ for all $x, y \in V$. We must show that $f(q x) = q f(x)$ for all $q \in \mathbb{Q}$ and $x \in V$.

\begin{enumerate}[label=\textbf{(\arabic*)}]
    \item \textbf{Integer Multiples:}
    First, $f(0_V) = f(0_V + 0_V) = f(0_V) + f(0_V) \implies f(0_V) = 0_W$.
    By induction, for any positive integer $n \in \mathbb{N}$:
    \[
    f(n x) = f(\underbrace{x + \dots + x}_{n \text{ times}}) = \underbrace{f(x) + \dots + f(x)}_{n \text{ times}} = n f(x).
    \]
    For negative integers, $0_W = f(0_V) = f(n x + (-n)x) = f(n x) + f((-n)x) = n f(x) + f((-n)x)$, so $f((-n)x) = -n f(x)$.
    Thus $f(m x) = m f(x)$ for all $m \in \mathbb{Z}$.

    \item \textbf{Rational Multiples:}
    Let $q = \frac{a}{b} \in \mathbb{Q}$ with $a, b \in \mathbb{Z}$ and $b \ne 0$.
    Apply the integer homogeneity result to the vector $y := \frac{1}{b} x \in V$:
    \[
    f(x) = f(b y) = b f(y) = b f\left(\frac{1}{b} x\right) \implies f\left(\frac{1}{b} x\right) = \frac{1}{b} f(x).
    \]
    Therefore,
    \[
    f(q x) = f\left(a \left(\frac{1}{b} x\right)\right) = a f\left(\frac{1}{b} x\right) = a \left(\frac{1}{b} f(x)\right) = \frac{a}{b} f(x) = q f(x).
    \]
\end{enumerate}
Thus $f$ is homogeneous under scalar multiplication by $\mathbb{Q}$, so $f \in \Hom_{\mathbb{Q}}(V, W)$.
\exinfo{Official Solution (Problem Sheet 11, Exercise 2). Establishes $\mathbb{Q}$-homogeneity for additive maps using induction and scaling.}
\end{exercisesolution}

% Official Solution (Problem Sheet 11, Exercise 3)
\begin{exercisesolution}[Solution to \cref{exc:polynomial_integral_eval_linear_map}]
\label{sol:polynomial_integral_eval_linear_map}
Write $T = (T_1, T_2)\transp \colon \mathbb{R}[x] \to \mathbb{R}^2$ where:
\[
T_1(p) = 3p(4) + 5p(6) + b p(1)p(2), \qquad T_2(p) = \int_{-1}^2 x^2 p(x) \, dx + c (p(0))^2.
\]
\begin{itemize}
    \item \textbf{If $b = c = 0$:} The components become:
    \[
    T_1(p) = 3p(4) + 5p(6), \qquad T_2(p) = \int_{-1}^2 x^2 p(x) \, dx.
    \]
    Evaluation at a point ($p \mapsto p(x_0)$) and definite integration against a fixed weight function are well-known linear functionals. Any linear combination of linear functionals is linear, so $T_1$ and $T_2$ are linear, making $T$ linear.

    \item \textbf{Conversely, suppose $T$ is linear:}
    Linearity requires $T(\alpha p) = \alpha T(p)$ for all $\alpha \in \mathbb{R}$ and $p \in \mathbb{R}[x]$.
    Consider the constant polynomial $p_0(x) = 1$:
    \[
    T_1(p_0) = 3(1) + 5(1) + b(1)(1) = 8 + b, \qquad T_1(2 p_0) = 3(2) + 5(2) + b(2)(2) = 16 + 4b.
    \]
    By linearity, $T_1(2 p_0) = 2 T_1(p_0) = 16 + 2b$. Equating the two yields $16 + 4b = 16 + 2b \implies 2b = 0 \implies b = 0$.

    Similarly for $T_2$, using $\int_{-1}^2 x^2 \, dx = \left[\frac{x^3}{3}\right]_{-1}^2 = \frac{8 - (-1)}{3} = 3$:
    \[
    T_2(p_0) = 3 + c(1)^2 = 3 + c, \qquad T_2(2 p_0) = 2(3) + c(2)^2 = 6 + 4c.
    \]
    By linearity, $T_2(2 p_0) = 2 T_2(p_0) = 6 + 2c$. Equating the two yields $6 + 4c = 6 + 2c \implies 2c = 0 \implies c = 0$.
\end{itemize}
Thus $T$ is linear if and only if $b = c = 0$.
\exinfo{Official Solution (Problem Sheet 11, Exercise 3). Identifies parameter values making integral and point-evaluation maps linear via test polynomials.}
\end{exercisesolution}

% Official Solution (Problem Sheet 11, Multiple Choice)
\begin{exercisesolution}[Solution to \cref{exc:identifying_linear_maps_mc}]
\label{sol:identifying_linear_maps_mc}
\begin{enumerate}[label=\textbf{(\arabic*)}]
    \item \textbf{Linear:} $f(x, y) = (x + y, 2x, 0)$ is given by matrix multiplication:
    \[
    f\begin{pmatrix} x \\ y \end{pmatrix} = \begin{pmatrix} 1 & 1 \\ 2 & 0 \\ 0 & 0 \end{pmatrix} \begin{pmatrix} x \\ y \end{pmatrix}.
    \]
    By \cref{ex:matrix_transformation}, every matrix transformation is linear.

    \item \textbf{Not Linear:} The second component contains the non-linear product $2xy$. For instance:
    \[
    f(1, 1) = (2, 2, 0) \implies 2 f(1, 1) = (4, 4, 0), \quad\text{but}\quad f(2 \cdot (1, 1)) = f(2, 2) = (4, 8, 0) \ne (4, 4, 0).
    \]
    Thus $f$ fails scalar homogeneity.

    \item \textbf{Linear:} The map is given by matrix multiplication by $\begin{pmatrix} \alpha & \beta & \gamma \\ \delta & \varepsilon & \eta \end{pmatrix} \in \M_{2 \times 3}(K)$, which is linear.
\end{enumerate}
\exinfo{Official Solution (Problem Sheet 11, Multiple Choice). Analyzes linearity of vector formulas and matrix operations.}
\end{exercisesolution}

% Official Solution (Problem Sheet 11, Multiple Choice)
\begin{exercisesolution}[Solution to \cref{exc:matrix_map_mc}]
\label{sol:matrix_map_mc}
The matrix $A = \begin{pmatrix} 2 & 1 \\ 1 & 0 \\ 0 & 2 \end{pmatrix}$ has size $3 \times 2$, so it defines a linear map $K^2 \to K^3$ via:
\[
A \begin{pmatrix} x \\ y \end{pmatrix} = \begin{pmatrix} 2 & 1 \\ 1 & 0 \\ 0 & 2 \end{pmatrix} \begin{pmatrix} x \\ y \end{pmatrix} = \begin{pmatrix} 2x + y \\ x \\ 2y \end{pmatrix}.
\]
Therefore, map \textbf{(1)}, $f \colon \mathbb{R}^2 \to \mathbb{R}^3, f(x, y) = (2x + y, x, 2y)\transp$, is the correct one.
\exinfo{Official Solution (Problem Sheet 11, Multiple Choice). Identifies the canonical linear transformation associated with a $3 \times 2$ matrix.}
\end{exercisesolution}

% Official Solution (Problem Sheet 11, Exercise 4)
\begin{exercisesolution}[Solution to \cref{exc:preimage_subspace_dimension_formula}]
\label{sol:preimage_subspace_dimension_formula}
\begin{enumerate}[label=\textbf{(\alph*)}]
    \item \textbf{Subspace Property:}
    Since $0_W \in W'$ and $f(0_V) = 0_W$, we have $0_V \in f^{-1}(W')$.
    Let $\alpha, \beta \in K$ and $v_1, v_2 \in f^{-1}(W')$, so $f(v_1), f(v_2) \in W'$.
    By linearity of $f$, $f(\alpha v_1 + \beta v_2) = \alpha f(v_1) + \beta f(v_2) \in W'$ because $W'$ is a linear subspace.
    Hence $\alpha v_1 + \beta v_2 \in f^{-1}(W')$, so $f^{-1}(W')$ is a linear subspace of $V$.

    \item \textbf{Dimension Formula:}
    Consider the restriction of $f$ to the subspace $f^{-1}(W')$:
    \[
    g := f|_{f^{-1}(W')} \colon f^{-1}(W') \to W.
    \]
    \begin{itemize}
        \item \textbf{Kernel of $g$:} $\ker(g) = \{v \in f^{-1}(W') \mid f(v) = 0\} = \{v \in V \mid f(v) = 0\} = \ker(f)$, since $\ker(f) = f^{-1}(\{0\}) \subseteq f^{-1}(W')$.
        \item \textbf{Image of $g$:} $\Img(g) = f(f^{-1}(W')) = \Img(f) \cap W'$.
    \end{itemize}
    Applying the Rank-Nullity Theorem (\cref{thm:rank_nullity}) to the linear map $g$:
    \[
    \dim f^{-1}(W') = \dim \ker(g) + \dim \Img(g) = \dim \ker(f) + \dim\big(\Img(f) \cap W'\big). \qedhere
    \]
\end{enumerate}
\exinfo{Official Solution (Problem Sheet 11, Exercise 4). Proves preimage of a subspace is a subspace and derives its dimension formula via Rank-Nullity on the restricted map.}
\end{exercisesolution}

% Official Solution (Problem Sheet 11, Exercise 5)
\begin{exercisesolution}[Solution to \cref{exc:isomorphisms_preserve_structures}]
\label{sol:isomorphisms_preserve_structures}
Let $T \colon V \to W$ be an isomorphism with linear inverse $T^{-1} \colon W \to V$.

\begin{enumerate}[label=\textbf{(\alph*)}]
    \item \textbf{Preservation of Linear Independence:}
    Let $S \subseteq V$ be linearly independent. Let $w_1, \dots, w_k \in T(S)$ be distinct vectors, so $w_i = T(v_i)$ for unique $v_i \in S$.
    Suppose $c_1 w_1 + \dots + c_k w_k = 0_W$.
    Applying the linearity of $T$:
    \[
    T(c_1 v_1 + \dots + c_k v_k) = c_1 T(v_1) + \dots + c_k T(v_k) = 0_W.
    \]
    Since $T$ is an isomorphism, it is injective, so $\ker(T) = \{0_V\}$. This implies $c_1 v_1 + \dots + c_k v_k = 0_V$.
    Since $S$ is linearly independent, $c_1 = \dots = c_k = 0$. Thus $T(S)$ is linearly independent in $W$.

    \item \textbf{Preservation of Spanning Sets:}
    Let $S \subseteq V$ span $V$, so $\Sp(S) = V$.
    For any $w \in W$, since $T$ is surjective, there exists $v \in V$ such that $T(v) = w$.
    Since $S$ spans $V$, we can write $v = \sum_{i=1}^m a_i v_i$ for some $v_i \in S$ and $a_i \in K$.
    Then:
    \[
    w = T(v) = T\left(\sum_{i=1}^m a_i v_i\right) = \sum_{i=1}^m a_i T(v_i) \in \Sp(T(S)).
    \]
    Thus $\Sp(T(S)) = W$.

    \item \textbf{Preservation of Bases:}
    A basis $\mathcal{B}$ of $V$ is both linearly independent and spans $V$.
    By parts \textbf{(a)} and \textbf{(b)}, $T(\mathcal{B})$ is linearly independent and spans $W$, hence $T(\mathcal{B})$ is a basis of $W$.
\end{enumerate}
\exinfo{Official Solution (Problem Sheet 11, Exercise 5). Proves isomorphisms preserve linear independence, spanning sets, and bases.}
\end{exercisesolution}

% Official Solution (Problem Sheet 11, Exercise 6)
\begin{exercisesolution}[Solution to \cref{exc:one_dim_endomorphism_canonical_iso}]
\label{sol:one_dim_endomorphism_canonical_iso}
\begin{enumerate}[label=\textbf{(\alph*)}]
    \item \textbf{Scalar Multiplication Property:}
    Since $\dim V = 1$, choose any non-zero vector $v_0 \in V$. Then $(v_0)$ is a basis of $V$.
    Since $T(v_0) \in V$, there exists a unique scalar $\lambda \in K$ such that $T(v_0) = \lambda v_0$.
    For any arbitrary vector $v \in V$, we have $v = c v_0$ for some $c \in K$. By linearity:
    \[
    T(v) = T(c v_0) = c T(v_0) = c (\lambda v_0) = \lambda (c v_0) = \lambda v.
    \]
    If $\lambda, \mu \in K$ both satisfy $T(v) = \lambda v = \mu v$ for all $v \in V$, choosing $v = v_0 \ne 0$ gives $(\lambda - \mu)v_0 = 0 \implies \lambda = \mu$, proving uniqueness.

    \item \textbf{Non-Canonical vs.~Canonical Isomorphisms:}
    \begin{itemize}
        \item An isomorphism $\psi \colon V \to K$ requires choosing a non-zero basis vector $v_0 \in V$ and setting $\psi(v_0) = 1_K$. If we replace $v_0$ by another basis vector $\alpha v_0$ (with $\alpha \in K^\times, \alpha \ne 1$), the corresponding isomorphism changes by a factor of $\alpha^{-1}$. Since there is no distinguished non-zero vector in $V$, there is no canonical isomorphism $V \cong K$.
        \item In contrast, every endomorphism $T \in \Hom_K(V, V)$ acts as $T = \lambda \id_V$ for a uniquely determined $\lambda \in K$. The map $\Phi \colon \Hom_K(V, V) \to K, T \mapsto \lambda$ is coordinate-free: for any non-zero vector $v \in V$, $\lambda$ is uniquely given by the ratio $T(v)/v$. One readily checks that $\Phi$ is an isomorphism of $K$-vector spaces (and even a ring isomorphism), providing a canonical identification $\Hom_K(V, V) \cong K$.
    \end{itemize}
\end{enumerate}
\exinfo{Official Solution (Problem Sheet 11, Exercise 6). Explains the canonical scalar isomorphism $\operatorname{End}(V) \cong K$ in dimension 1 versus non-canonical $V \cong K$.}
\end{exercisesolution}

% Official Solution (Problem Sheet 11, Exercise 7)
\begin{exercisesolution}[Solution to \cref{exc:functionals_same_kernel_proportional}]
\label{sol:functionals_same_kernel_proportional}
Let $T_1, T_2 \in \Hom_K(V, K)$ with $\ker(T_1) = \ker(T_2) =: U$.
\begin{itemize}
    \item \textbf{Case 1: $U = V$.}
    Then $T_1(v) = 0$ and $T_2(v) = 0$ for all $v \in V$, so $T_1 = 0 = 1 \cdot T_2$ (or $c T_2$ for any $c \in K$).

    \item \textbf{Case 2: $U \subsetneq V$.}
    Choose a vector $v_0 \in V \setminus U$. Since $v_0 \notin \ker(T_1)$ and $v_0 \notin \ker(T_2)$, we have $T_1(v_0) \ne 0$ and $T_2(v_0) \ne 0$.
    Define the scalar:
    \[
    c := \frac{T_1(v_0)}{T_2(v_0)} \in K^\times.
    \]
    Now let $v \in V$ be any vector. Define $\alpha := \frac{T_2(v)}{T_2(v_0)} \in K$, and consider $u := v - \alpha v_0 \in V$.
    Evaluating $T_2$ on $u$:
    \[
    T_2(u) = T_2(v - \alpha v_0) = T_2(v) - \alpha T_2(v_0) = T_2(v) - \frac{T_2(v)}{T_2(v_0)} T_2(v_0) = 0.
    \]
    Thus $u \in \ker(T_2) = U = \ker(T_1)$, which means $T_1(u) = 0$.
    Evaluating $T_1$ on $v = u + \alpha v_0$:
    \[
    T_1(v) = T_1(u) + \alpha T_1(v_0) = 0 + \frac{T_2(v)}{T_2(v_0)} T_1(v_0) = c T_2(v).
    \]
    Since this holds for all $v \in V$, we conclude $T_1 = c T_2$.
\end{itemize}
\exinfo{Official Solution (Problem Sheet 11, Exercise 7). Demonstrates that linear functionals with identical nullspaces are scalar multiples of each other.}
\end{exercisesolution}

% Official Solution (Problem Sheet 11, Exercise 8)
\begin{exercisesolution}[Solution to \cref{exc:projections_and_complements}]
\label{sol:projections_and_complements}
\begin{enumerate}[label=\textbf{(\alph*)}]
    \item \textbf{$V = \ker(P) \oplus \Img(P)$ for any projection $P$:}
    \begin{itemize}
        \item \textbf{Sum:} For any $v \in V$, write $v = (v - P(v)) + P(v)$.
        Applying $P$ to the first term and using $P^2 = P$:
        \[
        P(v - P(v)) = P(v) - P^2(v) = P(v) - P(v) = 0_V \implies v - P(v) \in \ker(P).
        \]
        Clearly $P(v) \in \Img(P)$. Thus $V = \ker(P) + \Img(P)$.
        \item \textbf{Trivial Intersection:} Let $v \in \ker(P) \cap \Img(P)$. Since $v \in \ker(P)$, $P(v) = 0$. Since $v \in \Img(P)$, $v = P(u)$ for some $u \in V$. Applying $P$ gives:
        \[
        v = P(u) = P^2(u) = P(P(u)) = P(v) = 0_V.
        \]
        Thus $\ker(P) \cap \Img(P) = \{0_V\}$.
    \end{itemize}
    Hence $V = \ker(P) \oplus \Img(P)$, so $\Img(P)$ is a complement of $\ker(P)$.

    \item \textbf{Existence and Uniqueness of Projection for a Direct Sum Decomposition:}
    Let $V = W_1 \oplus W_2$. Every $v \in V$ can be uniquely written as $v = w_1 + w_2$ with $w_1 \in W_1, w_2 \in W_2$.
    Define $P \colon V \to V$ by $P(v) := w_2$.
    \begin{itemize}
        \item \textbf{Linearity:} If $v = w_1 + w_2$ and $v' = w'_1 + w'_2$, then for any $\alpha, \beta \in K$, $\alpha v + \beta v' = (\alpha w_1 + \beta w'_1) + (\alpha w_2 + \beta w'_2)$. By uniqueness of direct sum representation, $P(\alpha v + \beta v') = \alpha w_2 + \beta w'_2 = \alpha P(v) + \beta P(v')$.
        \item \textbf{Idempotence:} For any $v \in V$, $P(v) = w_2 = 0 + w_2$, so $P(P(v)) = w_2 = P(v)$, which shows $P^2 = P$.
        \item \textbf{Kernel and Image:} $P(v) = 0 \iff w_2 = 0 \iff v = w_1 \in W_1$, so $\ker(P) = W_1$. Clearly $\Img(P) = W_2$.
        \item \textbf{Uniqueness:} Let $Q \colon V \to V$ be another projection with $\ker(Q) = W_1$ and $\Img(Q) = W_2$. For any $w_2 \in W_2 = \Img(Q)$, there exists $u \in V$ with $w_2 = Q(u)$, so $Q(w_2) = Q^2(u) = Q(u) = w_2$. For $w_1 \in W_1 = \ker(Q)$, $Q(w_1) = 0$. Thus for any $v = w_1 + w_2$, $Q(v) = Q(w_1) + Q(w_2) = 0 + w_2 = w_2 = P(v)$, proving $Q = P$.
    \end{itemize}
\end{enumerate}
\exinfo{Official Solution (Problem Sheet 11, Exercise 8). Establishes the 1-to-1 correspondence between direct sum decompositions and projection operators.}
\end{exercisesolution}

% Solution: Gemini / Official Hybrid (Problem Sheet 11, Exercise 9)
\begin{exercisesolution}[Solution to \cref{exc:invertibility_three_map_composition}]
\label{sol:invertibility_three_map_composition}
We are given $S, T, U \in \Hom(V, V)$ on a finite-dimensional space $V$ with $S \circ T \circ U = \id_V$.

\begin{itemize}
    \item Grouping the composition as $S \circ (T \circ U) = \id_V$: by \cref{cor:equal_dim_iso}, for endomorphisms on a finite-dimensional vector space, having a right inverse implies two-sided invertibility. Thus $S$ and $T \circ U$ are invertible with $S^{-1} = T \circ U$. This implies:
    \[
    (T \circ U) \circ S = \id_V \implies T \circ (U \circ S) = \id_V.
    \]
    \item Grouping the composition as $(S \circ T) \circ U = \id_V$: similarly, $U$ and $S \circ T$ are invertible with $U^{-1} = S \circ T$, which implies:
    \[
    U \circ (S \circ T) = \id_V \implies (U \circ S) \circ T = \id_V.
    \]
\end{itemize}
Since $(U \circ S) \circ T = \id_V$ and $T \circ (U \circ S) = \id_V$, the map $T$ is invertible with two-sided inverse:
\[
T^{-1} = U \circ S. \qedhere
\]
\exinfo{Hybrid Solution (Gemini 3.7 Flash / Official Problem Sheet 11, Exercise 9). Deduces invertibility of the middle endomorphism $T$ when $S \circ T \circ U = \id_V$ in finite dimensions.}
\end{exercisesolution}

% Official Solution (Problem Sheet 11, Exercise 10)
\begin{exercisesolution}[Solution to \cref{exc:endomorphism_properties_true_false}]
\label{sol:endomorphism_properties_true_false}
\begin{enumerate}[label=\textbf{(\alph*)}]
    \item \textbf{Statement (1) is TRUE:}
    Let $\mathcal{B}' = (v_1, \dots, v_k)$ be a basis for the subspace $V'$. Extend $\mathcal{B}'$ to a basis $\mathcal{B} = (v_1, \dots, v_k, v_{k+1}, \dots, v_n)$ of $V$.
    Define the linear map $\bar{f} \colon V \to V$ on the basis $\mathcal{B}$ by:
    \[
    \bar{f}(v_i) := f(v_i) \quad (1 \le i \le k), \qquad \bar{f}(v_j) := v_j \quad (k+1 \le j \le n).
    \]
    By definition, $\bar{f}|_{V'} = f$. To show $\bar{f}$ is an automorphism, we verify injectivity ($\ker(\bar{f}) = \{0_V\}$):
    Let $v = \sum_{i=1}^k a_i v_i + \sum_{j=k+1}^n b_j v_j \in \ker(\bar{f})$. Then:
    \[
    \bar{f}(v) = \underbrace{f\left(\sum_{i=1}^k a_i v_i\right)}_{\in V'} + \underbrace{\sum_{j=k+1}^n b_j v_j}_{\in W := \Sp(v_{k+1}, \dots, v_n)} = 0_V.
    \]
    Since $V = V' \oplus W$, the intersection $V' \cap W = \{0_V\}$, so both components must be zero:
    $\sum_{j=k+1}^n b_j v_j = 0 \implies b_j = 0$ for all $j$, and $f(\sum_{i=1}^k a_i v_i) = 0 \implies \sum_{i=1}^k a_i v_i = 0$ (since $f$ is an automorphism of $V'$), so $a_i = 0$ for all $i$.
    Thus $\ker(\bar{f}) = \{0_V\}$, so $\bar{f}$ is an automorphism.

    \item \textbf{Statement (2) is FALSE:}
    Consider the endomorphism $f \colon \mathbb{R}^2 \to \mathbb{R}^2$ given by $f(x, y) = (0, x)$.
    Then $\ker(f) = \{(x, y) \mid x = 0\} = \Sp((0, 1))$ and $\Img(f) = \{(0, x) \mid x \in \mathbb{R}\} = \Sp((0, 1))$.
    Since $\ker(f) = \Img(f)$, we have $\ker(f) \cap \Img(f) = \Sp((0, 1)) \ne \{0\}$, so $\Img(f)$ is not a complement of $\ker(f)$.

    \item \textbf{Statement (3) is TRUE:}
    For any linear map $T \colon \mathbb{R}^5 \to \mathbb{R}^5$, the Rank-Nullity Theorem (\cref{thm:rank_nullity}) states:
    \[
    \dim \ker(T) + \rank(T) = \dim \mathbb{R}^5 = 5.
    \]
    If $\rank(T) = \dim \ker(T) = k$, then $2k = 5 \implies k = \frac{5}{2}$. Since the dimension of a subspace must be an integer, no such linear map can exist.
\end{enumerate}
\exinfo{Official Solution (Problem Sheet 11, Exercise 10). Evaluates statements on automorphism extension, image-kernel complements, and parity of rank-nullity dimensions.}
\end{exercisesolution}
